\documentclass[11pt]{article}
\usepackage[margin=1in]{geometry}
\usepackage{amsmath,amssymb,amsthm}
\usepackage{enumitem}
\usepackage[T1]{fontenc}
\usepackage{lmodern}
\usepackage{hyperref}

\newtheorem{theorem}{Theorem}
\newtheorem{lemma}{Lemma}
\newtheorem{corollary}{Corollary}
\theoremstyle{definition}
\newtheorem{definition}{Definition}
\newtheorem{remark}{Remark}

\newcommand{\N}{\mathbb N}
\newcommand{\supp}{\operatorname{supp}}
\newcommand{\calG}{\mathcal G}
\newcommand{\calD}{\mathcal D}
\newcommand{\fp}{f_p}

\title{A logarithmic-democracy obstruction for the power-profile greedy question}
\author{Run \texttt{fa\_banach\_001}, lane 7}
\date{2026-06-21}

\begin{document}
\maketitle

\section*{Source question}

The source paper is Pablo M. Berna and Hung Viet Chu,
\emph{On Some Characterizations of Greedy-type Bases}, arXiv:2201.12029.
For a function \(f:\N\to\mathbb R\), and a finite set
\(A=\{n_1<\cdots<n_m\}\), write
\[
        1_{f,A}=\sum_{j=1}^m f(j)e_{n_j},
\]
and
\[
        \calD_m^f(x)=
        \inf\{ d(x,[1_{f,A}]): A\subset\N,\ |A|=m\}.
\]
A basis is \(f\)-greedy if there is \(C_f\) such that
\[
        \|x-\calG_m(x)\|\le C_f \calD_m^f(x)
        \qquad (x\in X,\ m\in\N).
        \tag{1}
\]
The first future-research question in the source asks for a characterization
of those functions \(f\) for which (1) implies ordinary greediness.  The source
proves this for regular \(f\), constructs non-greedy examples for sparse
\(f\), and explicitly points to
\[
        f(n)=n^{-p},\qquad p>0,
\]
as a class lying outside both results.

\section*{Result}

We do not settle the full power-profile problem.  The following partial result
shows, however, that the power profiles already force a strong quantitative
form of democracy.

\begin{theorem}\label{thm:main}
Let \(\mathcal B=(e_n)\) be a semi-normalized Schauder basis of a Banach
space \(X\).  Fix \(p>0\) and put \(\fp(n)=n^{-p}\).  If \(\mathcal B\) is
\(\fp\)-greedy with constant \(C_p\), then \(\mathcal B\) is suppression
unconditional and there is a constant \(L\), depending only on \(p\), \(C_p\),
and the semi-normalization constants of \(\mathcal B\), such that
\[
        \|1_A\|\le L\,\lceil \log_2(|A|+1)\rceil\, \|1_B\|
        \tag{2}
\]
whenever \(A,B\subset\N\) are finite and \(|A|=|B|\ge 1\).
\end{theorem}

Thus any counterexample to the full question for \(f(n)=n^{-p}\) must be
unconditional and can fail democracy by at most a logarithmic factor.

\section*{Proof intuition}

The source already observes that \(f\)-greediness implies unconditionality
when \(f\) has no zero values.  The new point is that the power profile cannot
hide large dyadic tail pieces.  If a block \(A\) has at least as many earlier
coordinates \(C\) before it, then
\[
        x=1_{\fp,C\cup A}+\eta 1_B,
        \qquad \fp(|C|+1)<\eta<\fp(|C|)
\]
has greedy set \(C\cup B\).  The greedy error is the power-weighted copy of
\(A\), while \(1_{\fp,C\cup A}\) itself is an admissible one-dimensional
competitor for \(\calD_m^{\fp}(x)\).  Since a dyadic tail of \(n^{-p}\) has
first/last ratio at most \(2^p\), every dyadic tail piece of \(A\) is
controlled by \(B\).  Summing the dyadic pieces costs only \(\log |A|\).

\section*{Proof}

Let
\[
        a=\inf_n \|e_n\|>0,\qquad b=\sup_n\|e_n\|<\infty.
\]
By Proposition 3.1 of the source paper, applied to the non-vanishing function
\(\fp\), the basis is suppression unconditional.  Let \(K\) be a suppression
unconditional constant.  We shall only use the two standard consequences
\[
        \|P_E x\|\le K\|x\|,
        \qquad
        \Big\|\sum_{n\in E} u_n e_n\Big\|
        \le K\Big\|\sum_{n\in E} v_n e_n\Big\|
        \quad\text{if } |u_n|\le |v_n|.
        \tag{3}
\]

\begin{lemma}[tail-block estimate]\label{lem:tail}
Let \(C,A,B\subset\N\) be pairwise disjoint finite sets with
\[
        |A|=|B|=r,\qquad |C|=s\ge r,\qquad \max C<\min A .
\]
Then
\[
        \|1_A\|\le C_p K\,2^p\,\|1_B\|.
        \tag{4}
\]
\end{lemma}

\begin{proof}
Write \(C\cup A=\{d_1<\cdots<d_{s+r}\}\).  Because \(\max C<\min A\), the
points of \(C\) are \(d_1,\ldots,d_s\), and the points of \(A\) are
\(d_{s+1},\ldots,d_{s+r}\).

Choose \(\eta\) with
\[
        \fp(s+1)<\eta<\fp(s).
\]
For
\[
        x=1_{\fp,C\cup A}+\eta 1_B,
\]
the natural greedy set of order \(s+r\) is \(C\cup B\): the \(C\)-coefficients
are at least \(\fp(s)\), the \(B\)-coefficients are \(\eta\), and the
\(A\)-coefficients are at most \(\fp(s+1)\).  Hence
\[
        x-\calG_{s+r}(x)
        =\sum_{j=1}^r \fp(s+j)e_{a_j},
\]
where \(A=\{a_1<\cdots<a_r\}\).  By (3),
\[
        \fp(s+r)\|1_A\|
        \le
        K\Big\|\sum_{j=1}^r \fp(s+j)e_{a_j}\Big\|
        =
        K\|x-\calG_{s+r}(x)\|.
        \tag{5}
\]
On the other hand, \(1_{\fp,C\cup A}\) is an allowed competitor in
\(\calD_{s+r}^{\fp}(x)\), so
\[
        \calD_{s+r}^{\fp}(x)
        \le \|x-1_{\fp,C\cup A}\|
        = \eta\|1_B\|.
        \tag{6}
\]
Combining (1), (5), and (6), and then letting
\(\eta\downarrow \fp(s+1)\), gives
\[
        \|1_A\|
        \le C_pK\,\frac{\fp(s+1)}{\fp(s+r)}\,\|1_B\|.
\]
Since \(s\ge r\),
\[
        \frac{\fp(s+1)}{\fp(s+r)}
        =
        \left(\frac{s+r}{s+1}\right)^p
        \le 2^p,
\]
which proves the lemma.
\end{proof}

\begin{lemma}[disjoint logarithmic democracy]\label{lem:disjoint}
There is a constant \(L_0\) such that
\[
        \|1_A\|\le L_0\,\lceil\log_2(|A|+1)\rceil\,\|1_B\|
        \tag{7}
\]
whenever \(A\cap B=\varnothing\) and \(|A|=|B|\ge 1\).
\end{lemma}

\begin{proof}
Let \(m=|A|=|B|\), and write \(A=\{a_1<\cdots<a_m\}\).  The singleton term is
controlled by
\[
        \|e_{a_1}\|\le b\le \frac{bK}{a}\|1_B\|.
        \tag{8}
\]
For \(j\ge1\), define the dyadic tail pieces
\[
        A_j=\{a_{2^{j-1}+1},\ldots,a_{\min(2^j,m)}\},
\]
omitting empty sets.  Let \(r_j=|A_j|\).  The set
\[
        C_j=\{a_1,\ldots,a_{2^{j-1}}\}
\]
has cardinal \(2^{j-1}\ge r_j\) and lies before \(A_j\).  Choose any
\(B_j\subset B\) with \(|B_j|=r_j\).  By Lemma \ref{lem:tail},
\[
        \|1_{A_j}\|\le C_pK2^p\|1_{B_j}\|
        \le C_pK^2 2^p\|1_B\|.
        \tag{9}
\]
There are at most \(\lceil\log_2(m+1)\rceil\) nonempty dyadic pieces, including
the singleton piece accounted for in (8).  Summing (8) and (9) gives (7), for
example with
\[
        L_0=\frac{bK}{a}+C_pK^2 2^p .
\]
\end{proof}

We now pass from disjoint sets to arbitrary sets.  Let \(A,B\subset\N\) have
\(|A|=|B|=m\).  Put \(I=A\cap B\), \(A_0=A\setminus B\), and
\(B_0=B\setminus A\).  Then \(|A_0|=|B_0|\), and the latter two sets are
disjoint.  If \(A_0=\varnothing\), there is nothing to prove.  Otherwise, by
suppression unconditionality and Lemma \ref{lem:disjoint},
\[
\begin{aligned}
        \|1_A\|
        &\le \|1_I\|+\|1_{A_0}\|  \\
        &\le K\|1_B\|
          + L_0\lceil\log_2(|A_0|+1)\rceil\,\|1_{B_0}\| \\
        &\le \bigl(K+K L_0\lceil\log_2(m+1)\rceil\bigr)\|1_B\|.
\end{aligned}
\]
Increasing the constant gives (2).  This proves Theorem \ref{thm:main}.

\section*{Verification notes and limitations}

The proof uses only the source paper's Proposition 3.1, the definition of
\(f\)-greediness, and elementary consequences of suppression unconditionality.
No computational check is involved.

The result is not a full solution of the source question for \(f(n)=n^{-p}\).
It proves that power-profile \(f\)-greedy bases have at most logarithmic
democracy distortion.  The remaining frontier is whether the logarithmic loss
can be removed, or whether there exists an \(\fp\)-greedy basis whose democracy
constants grow logarithmically.

\section*{Novelty check}

The run indexes were searched for arXiv:2201.12029, \(f\)-greedy,
\(\calD_m^f\), power-profile terminology, and \(1/n^p\).  Only an earlier
failed attempt for this same target was found.  Web searches on 2026-06-21 for
exact phrases around ``2201.12029'', ``f-greedy'', ``D\_m\^f'', and
``1/n\^p'' did not find a later resolution or a matching logarithmic-democracy
statement.

\section*{References}

\begin{enumerate}[label={[\arabic*]}]
\item P. M. Berna and H. V. Chu,
      \emph{On Some Characterizations of Greedy-type Bases},
      arXiv:2201.12029.
\item S. V. Konyagin and V. N. Temlyakov,
      A remark on greedy approximation in Banach spaces,
      East J. Approx. 5 (1999), 365--379.
\end{enumerate}

\end{document}
